\chapter{The Quantization Study}
\label{ch:quantization}

This chapter describes how the quantized models are built and how they are
compared with their full-precision parents. Every design choice below is either
inherited from the calibration of Section~\ref{sec:calibration} or settled by a
measurement reported in Section~\ref{sec:membrane-prep}; none is taken on the
test partition.

\section{Design principle: three separate arms}
\label{sec:grid}

The critical design choice is to \textbf{quantize the two targets separately
before quantizing them together}. Quantizing weights and membrane potential at
once would produce a single observation from which the contribution of each could
not be recovered; separating them turns the grid into a causal ablation and is
what makes the hypothesis of Chapter~\ref{ch:introduction} testable.

\begin{table}[htbp]
\centering
\resizebox{\linewidth}{!}{%
\begin{tabular}{llll}
\toprule
Arm & Weights & Membrane potential & Purpose \\
\midrule
FP32 reference & 32-bit float & 32-bit float & baseline for every comparison \\
Weight-only (W) & 8 / 4 / 2 bit & 32-bit float & isolates the synaptic (spatial) effect \\
Membrane-only (U) & 32-bit float & 8 / 4 / 2 bit & isolates the state (temporal) effect \\
Joint (W+U) & 8 / 4 / 2 bit & 8 / 4 / 2 bit & deployable configuration \\
\bottomrule
\end{tabular}}
\caption{The quantization grid: ten configurations, each at the three seeds of
the reference.}
\label{tab:grid}
\end{table}

Two quantities are deliberately left out of the grid. The membrane decay
$\beta = 0.75$ is already shift-implementable, having been constrained inside the
hyperparameter search (Section~\ref{sec:constrained-search}), and is not
re-quantized, which would be a third, unplanned intervention. The biases are kept
in full precision: there are 244 of them, under $0.1\,\%$ of the parameters,
and on integer hardware they are added once to a wide accumulator.

The two arms have \emph{opposite} sensitivities to the resolution choice, which
independently justifies treating them as separate ablations. Because $270\,240$
of the $295\,332$ parameters are convolutional and therefore
resolution-independent, the weight arm achieves nearly the same absolute memory
reduction at any resolution, whereas the membrane state falls twelvefold between
224 and 64 pixels.

\section{Quantization-aware training from the parent}
\label{sec:qat}

Quantization-aware training is implemented by fake quantization: in the forward
pass each quantized tensor is rounded to its grid and clamped to the grid's
bounds, and in the backward pass the rounding is replaced by the identity, the
straight-through estimator~\cite{bengio2013ste}.

\paragraph{Every quantized model is fine-tuned from the full-precision
checkpoint of the same seed.} This is the constraint that determines the rest of
the design. The numerical floor of the calibration ($\mathrm{PCC} = 0.975$) holds
for comparisons between a model and its own parent, where no channel permutation
exists; a quantized model trained from a random initialization would be an
independent solution, for which the relevant level is $0.153$ and almost no
quantization effect would be distinguishable from initialization variance. The
fine-tuning is therefore short and at reduced learning rate, adapting the parent
to its grid rather than finding a new solution. The relative $L_2$ distance
between the quantized weights and those of the parent is reported per layer, as
a check that each quantized model stayed in its parent's neighbourhood.

\todo{State the fine-tuning schedule actually used (number of epochs, learning
rate as a fraction of $1.047 \times 10^{-4}$, early-stopping rule) once fixed in
\texttt{train.py}.}

\paragraph{Execution order and sanity gate.} All 8-bit configurations are run
first, then 4-bit, then 2-bit. Before the grid is launched, the joint 8-bit model
of one seed must match its parent to within the spread of the full-precision
reference, four samples out of $1\,435$. At 8 bits the weight step is of the
order of the perturbation that left the spatial correlation at $0.975$ in the
calibration, so an appreciable loss of accuracy at that bit-width would indicate
a defect of the implementation, not a physical effect.

\paragraph{Fallback.} Should quantization-aware training prove unstable at the
lowest bit-widths, those configurations are produced by post-training
quantization of the parent instead, the procedure used by Rabbi et
al.~\cite{rabbi2026}; the comparison between the two procedures is then itself
reported.

\section{Weight quantizer}
\label{sec:weight-quantizer}

Weights are quantized on a \emph{symmetric signed} grid with
$q_{\max} = 2^{b-1} - 1$ levels on each side of zero, so that zero is exactly
representable: 255 levels at 8 bits, 15 at 4 bits and 3 at 2 bits (the values
$-\Delta, 0, +\Delta$, a ternary network). The quantizer is applied to every
convolutional and linear weight tensor, including the readout.

\textbf{The scale is computed per output channel.} The choice was made on the
measured weight distributions of the reference checkpoint
(Table~\ref{tab:weight-ranges}). In the deeper convolutions the largest weight
of the weakest channel is less than half the largest weight of the whole tensor,
so a per-tensor scale would leave that channel using fewer than half of the
levels: about one effective bit lost, which at 2 bits is the difference between
a ternary channel and a channel quantized entirely to zero.

\begin{table}[htbp]
\centering
\begin{tabular}{lcccc}
\toprule
Layer & $\max|w|$, tensor & $\max|w|$, weakest channel & Ratio & Bits lost per-tensor \\
\midrule
conv1  & 0.144 & 0.115 & 1.25 & 0.33 \\
conv2  & 0.126 & 0.060 & 2.11 & 1.08 \\
conv3  & 0.113 & 0.054 & 2.11 & 1.08 \\
conv4  & 0.099 & 0.044 & 2.26 & 1.18 \\
linear & 0.214 & 0.146 & 1.46 & 0.55 \\
\bottomrule
\end{tabular}
\caption{Weight ranges of the reference checkpoint (seed 1234). ``Bits lost'' is
$\log_2$ of the ratio: the resolution that the weakest channel would forgo under
a single per-tensor scale.}
\label{tab:weight-ranges}
\end{table}

The global 8-bit step used in the calibration ($0.00168$, from the largest
absolute weight of the network, $0.214$) is therefore an upper bound on the
per-channel steps actually applied.

\section{Membrane quantizer}
\label{sec:membrane-quantizer}

The membrane potential of the four hidden LIF layers is quantized at every time
step. The quantizer acts on the potential returned by each layer after the spike
decision, so the quantized value is the \emph{state} carried into the next step:
what is stored, and therefore what costs memory, is what is quantized. The
output layer, four neurons, is left in full precision; its state is negligible in
memory and quantizing it would act directly on the logits.

A grid for the membrane cannot be designed without knowing the range of the
quantity, and unlike a weight, the potential has no range until the network is
run. Its distribution was therefore measured over the whole validation partition
(Section~\ref{sec:membrane-prep}), and the clipping range was chosen by testing
candidate ranges on the full-precision network before any quantization.

\textbf{The membrane grid is unsigned over $[0, 2]$}, twice the firing threshold
$U_{\text{thr}} = 1$. The step is $2/(2^b - 1)$:

\begin{table}[htbp]
\centering
\begin{tabular}{cccc}
\toprule
Bits & Levels & Step & Step / threshold \\
\midrule
8 & 256 & 0.0078 & 0.8\,\% \\
4 & 16  & 0.1333 & 13.3\,\% \\
2 & 4   & 0.6667 & 66.7\,\% \\
\bottomrule
\end{tabular}
\caption{The membrane grid at each bit-width. At 2 bits the representable
potentials are $0$, $2/3$, $4/3$ and $2$: a single level lies strictly between
rest and threshold.}
\label{tab:membrane-grid}
\end{table}

Table~\ref{tab:membrane-grid} contains a quantitative prediction. At 8 bits the
rounding error is below one percent of the threshold; at 4 bits it is an eighth
of it; at 2 bits the error is of the same order as the threshold itself, and a
neuron's state can encode only whether it is at rest, below threshold, or above
it. This is the mechanism behind the observation of Smith et
al.~\cite{smith2026} that membrane quantization needs at least four bits, and the
grid tests it directly.

The unsigned scheme is a consequence of the clipping study rather than an
assumption: the network turns out to make no use of negative potentials, so a
signed grid would spend half of its levels on values that carry no information,
the equivalent of losing one bit at every bit-width.

\section{Joint arm}

The joint arm applies both quantizers at once, with the per-channel weight scales
of Section~\ref{sec:weight-quantizer} and the fixed membrane grid of
Section~\ref{sec:membrane-quantizer}, at the same bit-width.

\todo{Decide whether the joint arm keeps independent scales (as currently
implemented in \texttt{model/quantization.py}) or shares a single scale between
weights and membrane in the style of MINT~\cite{yin2024mint}, as stated in the
progress report. If scales stay independent, remove ``MINT-style'' everywhere
and justify the choice here; if shared, describe how the shared scale is set.}

\section{Reading the results}
\label{sec:reading}

For every quantized model, SAM maps are recomputed on the same fifteen validation
images per class and compared with those of \textbf{the full-precision parent of
the same seed}, using the metrics of Section~\ref{sec:metrics}: spatial PCC,
top-20\,\% IoU, the temporal centre of mass, and the per-timestep correlation
curve. Structural similarity~\cite{wang2004ssim} is added for comparability with
Rabbi et al.; unlike the other metrics it is not scale-invariant, and the maps
are normalized before it is computed.

Three further measurements characterize behaviour rather than explanations:
\begin{itemize}
  \item \textbf{cross-faithfulness}: the deletion metric applied to the quantized
        model with its \emph{own} map and with the map of its \emph{parent}. The
        second is the operational question, since in practice one is tempted to
        explain a deployed model with the maps of the original;
  \item the \textbf{Earth Mover's Distance} between the per-neuron firing-rate
        distributions of quantized model and parent, for direct comparison with
        Smith et al.~\cite{smith2026};
  \item the \textbf{diagnostics} of Section~\ref{sec:diagnostics} --- dead and
        saturated fractions, CV-ISI --- since quantization can induce dead or
        saturated neurons, which is a candidate mechanism for \emph{why} an
        explanation changes.
\end{itemize}

Five rules, all established in Chapter~\ref{ch:results} before the grid was
run, govern how those numbers are read.

\begin{enumerate}
  \item \textbf{Below the numerical floor is noise.} A divergence smaller than
        that produced by weight noise of one 8-bit step ($\mathrm{PCC}$ $0.975$,
        IoU $0.704$, CoM difference $0.029$ steps) is not a result.
  \item \textbf{IoU is an early indicator, PCC a graded measure.} IoU has a
        dynamic range of $2.6\times$ its floor and saturates early; PCC has
        $34\times$.
  \item \textbf{Everything is reported per class.} The sensitivity ordering
        Violin\_Mode $>$ Scattered\_Light $>$ Blip $>$ Extremely\_Loud recurs
        under every intervention tested, so a pooled figure would average a class
        that responds strongly with one that does not respond at all.
  \item \textbf{A flat result on Extremely\_Loud is uninformative.} That class
        did not register any perturbation tested; absence of effect and inability
        to detect one cannot be told apart.
  \item \textbf{Cross-faithfulness is interpretable on Blip and Extremely\_Loud
        only}, the two classes on which the deletion metric separates SAM from a
        random baseline. On the other two there is no signal for quantization to
        degrade.
\end{enumerate}

\paragraph{The test of the hypothesis.} The hypothesis predicts an
\emph{ordering}: the membrane arm should displace the temporal centre of mass
beyond its floor at a higher bit-width than the one at which the weight arm
lowers the spatial correlation beyond its floor. The reverse order refutes it.

\section{Stratified analysis}
\label{sec:stratified-method}

An aggregate divergence per configuration admits only a yes-or-no answer.
Joining the per-sample outputs to the physical descriptors allows the result to
be reported as a relation instead.

The hypothesis is that divergence increases as event strength decreases, with a
mechanism specific to spiking networks: quantizing the membrane potential flips a
neuron's spike decision only where the potential lies within the quantization
error of the threshold, and weak events produce potentials that sit close to
threshold rather than overshooting it. Because a flipped spike is a discrete
event that propagates both to the next layer and forward in time through the
decay and reset, and because SAM is constructed from spike trains, the flip
changes the object the explanation is made of rather than merely perturbing it.

The prediction discriminates between arms: membrane quantization applies a fixed
absolute error and should show a steep dependence on event strength, whereas
weight quantization perturbs the input current approximately in proportion to the
signal and should be flatter.

Two preconditions were verified. SNR does translate into image intensity, the
Gravity Spy colour scale being fixed rather than auto-scaled per image. And SNR
is very nearly a proxy for class --- the top global SNR quartile is 92.8\,\%
Extremely\_Loud --- so stratification uses \textbf{terciles computed within each
class}, or any trend measured would be a class effect wearing a strength label.

Finally, the \emph{disagreement subset} --- validation images classified
correctly by the parent and incorrectly after quantization --- is characterized
by its physical descriptors, to ask whether quantization fails on a particular
kind of event.

\section{Energy and memory}
\label{sec:energy-method}

Synaptic operations are counted per layer as
\begin{equation}
\mathrm{SOP}_{\ell} = f_{\ell-1} \cdot T \cdot \mathrm{MAC}_{\ell},
\end{equation}
where $f_{\ell-1}$ is the mean firing rate of the preceding hidden layer and
$\mathrm{MAC}_{\ell}$ the multiply-accumulate count of the equivalent dense
layer, and multiplied by the per-operation costs of~\cite{horowitz2014}. The
rates are the per-hidden-layer figures measured on the converged checkpoint of
each configuration, never the output-layer quantity logged as
\texttt{avg\_firing\_rate}.

\textbf{One layer is asymmetric.} Under direct encoding the first convolution
receives analog pixel values rather than spikes, so its operations are
multiply-accumulates. Because the input is identical at every time step, the
current it produces is constant in time and is computed once rather than $T$
times. This is a standard implementation of constant encoding, but stating it
matters: recomputing that layer at every step would make the spiking network
more expensive than its convolutional equivalent (Section~\ref{sec:energy-fp32}).

Memory is the number of bits needed to store the weights plus those needed to
store the membrane state of the $94\,080$ hidden neurons at the working
resolution. For integer configurations the 8-bit energy figures of Horowitz are
used; at 4 and 2 bits they are a conservative upper bound, since narrower
operations are not tabulated.

\section{Test-set evaluation}

The test partition is opened once, after every design decision has been taken,
on the full-precision reference and on the quantized configurations selected for
the final comparison. Accuracy is reported with a 95\,\% Wilson interval and per
class, and is computed by two routes --- an aggregate over batches and a
per-sample classification --- asserted equal, since the per-sample output feeds
the stratified analysis.
